\documentclass[track=junior,level=1,teacher]{kidscode}
\newcommand{\ostar}{\tikz[baseline=-0.6ex]\node[star,star points=5,star point ratio=2.3,draw=black,line width=0.8pt,fill=white,minimum size=0.6cm,inner sep=0pt]{};}
%% a script with a heading: \begin{refscript}{Sprite}{what it does} blocks \end{refscript}
\newenvironment{refscript}[2]{\par\needspace{8\baselineskip}\noindent\textbf{\color{kcaccent}#1:} #2\par\smallskip\begin{center}\begin{scratch}[scale=0.8]}{\end{scratch}\end{center}}
\begin{document}
\kctitle{Final Project Teacher Guide}{My Animated Greeting Card: milestones, stars, reference card and showcase}
\tableofcontents

%% =====================================================================
\section{The project in one page}
Each child builds an animated greeting card in Scratch for someone special, planned first on a
three-box storyboard (Start, Middle, End). The card must use five things: \textbf{the green flag to
start, at least one repeat (an animation or a dance), one click or key event, a sound, and two
sprites or a backdrop change}. In week 10 each child shows the card to the family, says what one
script does, and receives a certificate.

\begin{center}\small\renewcommand{\arraystretch}{1.35}
\begin{tabularx}{\linewidth}{@{}>{\bfseries\RaggedRight}p{3.0cm}>{\bfseries\RaggedRight}p{3.0cm}X@{}}
\toprule
When & Milestone & What the children do\\
\midrule
Week 9, warm-up (about 10 min) & M1 Storyboard & Watch your reference card, choose who the card is for, draw Start, Middle and End, check the must-haves on the plan.\\
Week 9, Make it (about 20 min) & M2 Card works & Build one storyboard box at a time as Driver and Helper; save with a nickname file name.\\
Week 9, Share & (check) & A partner checks the card against the storyboard.\\
Week 10 (about 30 min) & M3 Showcase & Polish, practice the 1-minute show-and-tell, family gallery, sticker feedback, certificates.\\
\bottomrule
\end{tabularx}
\end{center}

\begin{materials}
  \item Print one \textbf{Project Brief} per child (the storyboard page is the one they draw on).
  \item Before week 9: build the reference card yourself (section \enquote{The reference card} below,
        about 20 minutes) and save it as \texttt{Happy-Birthday-card} in \texttt{Scratch Projects}.
  \item For the core path: a \texttt{card-start} project with the backdrop and the first sprite already in place (see below).
  \item Week 10: the week-10 show-me checklist, the sticker feedback cards (last page), stickers,
        one certificate per child (\texttt{extras/\allowbreak{}certificate}), and the progress trackers.
\end{materials}

\begin{teachernote}[Why this project]
The card pulls every Level 1 goal together: exact steps in order (J1.G1), sequences and repeat
(J1.G2), events (J1.G3), fixing bugs (J1.G4), a storyboard (J1.G5), Driver and Helper and explaining
(J1.G6), and nicknames and credit (J1.G7). A card has no \enquote{losing}, every child can finish
something they are proud of, and it is made for a real person, which is a powerful reason to finish.
\end{teachernote}

%% =====================================================================
\section{Milestone walkthrough}

\subsection{M1 Storyboard (week 9 warm-up, about 10 minutes)}
\begin{teacherscript}
\enquote{Today you start your own project: a greeting card that moves and talks, made in Scratch for
someone special. Let me show you mine.} (Click the green flag on the reference card, then click the
balloon, then press the space bar.) \enquote{What did you see first? What happened when I clicked?
Before we build, we draw a plan, called a storyboard: what happens at the start, in the middle and at the end.}
\end{teacherscript}
\begin{enumerate}[itemsep=2pt]
  \item Read the mission and the must-haves aloud together. Ask: \enquote{Where was the repeat in my
        card? Where was the click?} (The cat's dance; the balloon.)
  \item Show the idea menu. Children choose who the card is for (first name or nickname only) and
        what kind of card. Offer two choices if a child cannot decide.
  \item They draw the three boxes: \textbf{Start} (what you see when the flag is clicked),
        \textbf{Middle} (the dance or animation), \textbf{End} (the click or key surprise).
        Stick figures and a few words are perfect.
  \item Point at each must-have icon on the page and ask \enquote{Where is this in your plan?}
        They check the boxes.
\end{enumerate}
\begin{formativecheck}[M1 is done when]
Three boxes have a drawing, the card has a name for who it is for, and the child can point to where
the repeat, the click or key and the sound will be. Quick question: \enquote{What happens when someone
clicks your card?}
\end{formativecheck}

\subsection{M2 Card works (week 9 Make it, about 20 minutes)}
Driver and Helper swap on the timer as usual: in a pair, they build one card together first (one
child's storyboard), then the other's, or each makes their own card on their own turn. With one child,
you are the Helper. Build in this order and click the flag after each step:
\begin{enumerate}[itemsep=2pt]
  \item \textbf{Start box:} backdrop and first sprite (core path: already there), a green-flag script
        with a \enquote{say} block. Click the flag.
  \item \textbf{Middle box:} the repeat with \enquote{next costume} and a wait. Click the flag.
  \item \textbf{End box:} a second sprite with \enquote{when this sprite clicked} or \enquote{when space key pressed}, and a sound.
  \item \textbf{Save:} \menu{File > Save to your computer}, into \texttt{Scratch Projects}, named with
        the nickname and \enquote{card}, for example \texttt{Tiger-card}.
\end{enumerate}
\begin{teacherscript}[Ask while you walk around]
\enquote{Which box of your storyboard are you building now?} \enquote{Show me the script that starts
the dance.} \enquote{What will happen when I click this sprite? Let's test it.}
\end{teacherscript}
\begin{formativecheck}[M2 is done when]
The card starts with the green flag, has all five must-haves, matches the storyboard (more or less),
and is saved. In Share, a partner checks the card against the storyboard: \enquote{Is the Start box there?
The Middle? The End?}
\end{formativecheck}

\subsection{M3 Showcase (week 10, about 30 minutes)}
\begin{enumerate}[itemsep=2pt]
  \item \textbf{Polish (8 minutes):} open the card (\menu{File > Load from your computer}). One small
        change only: a different word, a new sound, or a color. Save.
  \item \textbf{Practice (5 minutes):} each child practices the 1-minute show-and-tell with you or a partner.
  \item \textbf{Gallery (15 minutes):} see the showcase plan below.
  \item \textbf{Celebrate:} sticker feedback cards, the stars, the certificate.
\end{enumerate}
Fill in the week-10 show-me checklist while you watch: most rows show up during the polish and the
show-and-tell.

%% =====================================================================
\section{Unblocking common problems}
\begin{center}\small\renewcommand{\arraystretch}{1.3}
\begin{longtable}{@{}>{\RaggedRight}p{5.2cm}>{\RaggedRight\arraybackslash}p{\dimexpr\linewidth-5.2cm-2\tabcolsep}@{}}
\toprule
\textbf{Problem} & \textbf{What to do}\\
\midrule\endhead
\bottomrule\endfoot
Cannot choose an idea. & Offer two from the idea menu. \enquote{Birthday or thank-you?} The person comes first: \enquote{Who would love a card?}\\
The storyboard is huge. & Circle one thing in each box to build first. The rest is for later, if there is time.\\
The dance is too fast to see. & Put \enquote{wait 1 seconds} inside the repeat (the week-3 bug).\\
The dance happens only once. & \enquote{next costume} is under the repeat, not inside it. Drag it in until the orange mouth closes around it.\\
The balloon pops but never comes back. & Add \enquote{when green flag clicked} and \enquote{show} to the balloon. Hidden sprites stay hidden until something shows them.\\
Clicking the sprite does nothing. & The script is under another sprite. Click the sprite below the stage and check its code area.\\
The key does nothing. & Check the key in the dropdown. Click once on the stage so Scratch is listening, then press the key.\\
Two sprites talk at the same time. & Add a wait before the second sprite talks (the week-7 bug).\\
The backdrop does not change. & Check the new backdrop is in the \menu{Backdrops} tab (\menu{Choose a Backdrop}, bottom right) and its name matches the dropdown.\\
No sound. & Check the sprite has the sound in its \menu{Sounds} tab, and the Mac volume. See the Teacher Guide troubleshooting.\\
A real name is on the card. & Change it together to a nickname or a first name only. Explain why kindly.\\
\enquote{I want to add more!} & Great: first all five must-haves and a save, then extras one at a time.\\
The project is lost. & \menu{File > Load from your computer}, look in \texttt{Scratch Projects} and Downloads. Save every 10 minutes.\\
Tears or frustration. & Wiggle break. Then click the flag and find one thing that works: \enquote{Look, your cat already says Happy Birthday!}\\
\end{longtable}
\end{center}

%% =====================================================================
\section{The stars and how to use them kindly}
The brief has four rows, each with one to three stars. The stars are for the child, not a score to
report: there is no total, and children are never compared.

\begin{rubric}[My stars (adult version)]
\small\renewcommand{\arraystretch}{1.35}
\begin{tabularx}{\linewidth}{@{}>{\bfseries\RaggedRight}p{3.4cm}*{3}{>{\RaggedRight\arraybackslash}X}@{}}
 & \ostar & \ostar\,\ostar & \ostar\,\ostar\,\ostar\\
\midrule
It works & The card starts, but parts do not run. & Most scripts run when the flag, click or key is used. & Every script runs, every time, from a fresh green-flag click.\\ \midrule
It has every must-have & 1 or 2 of the 5. & 3 or 4 of the 5. & All 5: green flag, repeat, click or key, sound, two sprites or a backdrop change.\\ \midrule
I can say what one script does & Says what the card does overall. & Points to a script and names some blocks. & Picks one script and says what each block does, in order.\\ \midrule
I planned it on my storyboard & Drew a storyboard. & The card has some of the plan. & The card follows Start, Middle and End (changes along the way are fine).\\
\end{tabularx}
\end{rubric}

\begin{tip}[Using the stars kindly]
\begin{itemize}[itemsep=2pt]
  \item \textbf{The child colors first.} Read each row aloud and let them choose; then say what you saw.
        If you see more stars than they colored, say so. If you see fewer, ask \enquote{Can you show me?}
  \item \textbf{Two moments:} once at the end of week 9 (\enquote{what is left to do?}) and once after
        the showcase (\enquote{look how far you came}).
  \item \textbf{One star is a start.} Pick the one row to work on next and help with it right away.
  \item \textbf{Celebrate the plan and the explaining} as much as the working card.
\end{itemize}
\end{tip}

%% =====================================================================
\section{The reference card: \enquote{Happy Birthday!}}
Build this yourself before week 9 so you can show it, then keep it on \textbf{your} Mac. It is one
complete answer: children make their own card, not a copy. It uses only Level 1 blocks, whole
numbers, named turns and sprites dragged into place.

\subsection{Setting it up (no blocks yet)}
\begin{enumerate}[itemsep=2pt]
  \item New project (\menu{File > New}). The Cat is already there.
  \item Stage: click \menu{Choose a Backdrop} (bottom right) and pick \textbf{Party}. Keep the white
        \textbf{backdrop1} too: the card starts on it.
  \item Sprites: click \menu{Choose a Sprite} and add \textbf{Balloon1} and \textbf{Ballerina}.
        Drag the Cat to the left, the Ballerina to the right and the balloon to the top of the stage.
  \item Save as \texttt{Happy-Birthday-card} in \texttt{Scratch Projects}.
\end{enumerate}

\subsection{The scripts, block by block}
\begin{refscript}{Stage}{When the flag is clicked, the stage is plain for 1 second, then the party lights come on (the backdrop change).}
\blockinit{when \greenflag clicked}
\blocklook{switch backdrop to \selectmenu{backdrop1}}
\blockcontrol{wait \ovalnum{1} seconds}
\blocklook{switch backdrop to \selectmenu{Party}}
\end{refscript}

\begin{refscript}{Cat}{The cat says Happy Birthday!, meows, then dances: the repeat changes its costume 4 times, with a wait so we can see each step.}
\blockinit{when \greenflag clicked}
\blocklook{say \ovalnum{Happy Birthday!} for \ovalnum{2} seconds}
\blocksound{start sound \selectmenu{Meow}}
\blockrepeat{repeat \ovalnum{4}}{
  \blocklook{next costume}
  \blockcontrol{wait \ovalnum{1} seconds}
}
\end{refscript}

\begin{refscript}{Balloon1, script 1}{When the flag is clicked, the balloon shows (so it comes back after it pops).}
\blockinit{when \greenflag clicked}
\blocklook{show}
\end{refscript}

\begin{refscript}{Balloon1, script 2}{When someone clicks the balloon, it pops: the Pop sound plays and the balloon hides (the click event and the sound).}
\blockinit{when this sprite clicked}
\blocksound{start sound \selectmenu{Pop}}
\blocklook{hide}
\end{refscript}

\begin{refscript}{Ballerina}{When the space key is pressed, the Ballerina spins: four quarter turns make a whole spin, with a new costume each time.}
\blockinit{when \selectmenu{space} key pressed}
\blockrepeat{repeat \ovalnum{4}}{
  \blockmove{turn \turnright{} \ovalnum{90} degrees}
  \blocklook{next costume}
  \blockcontrol{wait \ovalnum{1} seconds}
}
\end{refscript}

\begin{center}\small\renewcommand{\arraystretch}{1.3}
\begin{tabularx}{\linewidth}{@{}>{\bfseries}p{4.6cm}X@{}}
\toprule
Must-have & Where it is in the reference card\\
\midrule
Green flag start & The Stage, the Cat and the balloon all start on the green flag.\\
A repeat & The Cat's dance (and the Ballerina's spin).\\
A click or key event & Click the balloon; press the space bar for the Ballerina.\\
A sound & Meow from the Cat, Pop from the balloon.\\
Two sprites or a backdrop change & Three sprites, and the backdrop changes from backdrop1 to Party.\\
\bottomrule
\end{tabularx}
\end{center}

\begin{tip}[Questions to ask about the reference card]
\enquote{How do we make the cat dance for longer?} (Change 4 to a bigger number.)
\enquote{What if we take the wait out of the dance?} (Too fast to see: try it!)
\enquote{Why does the balloon need a show block?} (Without it, once popped it stays hidden.)
\enquote{How many quarter turns make the Ballerina spin all the way round?} (Four.)
\enquote{Whose pictures and sounds are these?} (The Scratch library: we say so at the showcase.)
\end{tip}

\subsection{The core-path starting card}
For a child who needs the core path, make a \texttt{card-start} project before week 9: a new project,
choose the backdrop the child picked on the storyboard (or \textbf{Party}), keep the Cat, and save it
as \texttt{card-start} in \texttt{Scratch Projects}. In week 9 the child opens it and saves it straight
away with their own name, such as \texttt{Tiger-card}. They then add one must-have at a time.

\subsection{The stretch path: a pen frame}
The brief shows a Star sprite that draws a square frame with the pen and stamps itself at each
corner: pen down, then repeat 4 of stamp, move 100 steps and a quarter turn. Children drag the Star
to a corner of the stage first. Because a square ends where it started, the frame draws in the same
place every time the flag is clicked; \enquote{erase all} clears the old one.

%% =====================================================================
\section{Showcase plan (week 10)}

\subsection{The family gallery}
The showcase is at home, for family only. Each card gets a \enquote{stand}: a Mac with the card open and the stage in full screen (the button above the stage), the
storyboard beside it, and a folded card with the project name and the maker's \textbf{nickname}.
Ten to fifteen minutes is plenty with four children.

\begin{center}\small\renewcommand{\arraystretch}{1.3}
\begin{tabularx}{\linewidth}{@{}>{\bfseries}p{3.0cm}X@{}}
\toprule
Role & Job\\
\midrule
Card maker & Clicks the flag for each visitor and gives the 1-minute show-and-tell.\\
Helper & (A partner or a sibling.) Hands out feedback cards and stickers, and says \enquote{Click the balloon!}\\
Visitors & Watch, try the click or key, ask one question, put stickers on a feedback card.\\
Host (you) & Welcomes the family, keeps time, makes sure every child gets the same attention, reads the safety reminder.\\
\bottomrule
\end{tabularx}
\end{center}

\subsection{The 1-minute show-and-tell script}
Write it on a sticky note for each child. Practice once before the family arrives. Children who do not
read yet can point to pictures you draw next to each line.
\begin{enumerate}[itemsep=2pt]
  \item \enquote{Hi! My card is for \underline{\hspace{3cm}}.}
  \item \enquote{Watch: I click the green flag.} (Click.)
  \item \enquote{Now you click \underline{\hspace{2.5cm}}} or \enquote{press \underline{\hspace{2cm}}.}
  \item Point to one script: \enquote{This script makes \underline{\hspace{3cm}}. First it \underline{\hspace{2.5cm}}, then it \underline{\hspace{2.5cm}}.}
  \item \enquote{The pictures and sounds are from the Scratch library.}
  \item \enquote{Do you have a question for me?}
\end{enumerate}

\subsection{Certificates and closing}
Read the feedback cards aloud to each child, then hand out the certificate (write the nickname or
first name by hand). Close with a round of \enquote{one thing I am proud of, one thing I want to make
next}, and color the last stars on the progress tracker.

\begin{safety}[Showcase safety]
\begin{itemize}[itemsep=1pt]
  \item Family only, at home, on the screen. Nothing is uploaded, posted or emailed.
  \item Nicknames or first names only on the screen and on the stand cards.
  \item Photos only if the child agrees, with no names on the screen, and kept in the family.
  \item Credit: every card says the pictures and sounds come from the Scratch library.
\end{itemize}
\end{safety}

\begin{teachernote}[Fitting it into the 45-minute session]
0--10 polish and save (warm-up); 10--15 practice the show-and-tell (show me); 15--33 gallery and
show-and-tell, week-10 checklist (make it); 33--35 wiggle break; 35--45 feedback cards, stars,
certificates (share).
\end{teachernote}

\clearpage
\section{Sticker feedback cards (print and cut)}
Visitors put a sticker next to what they liked, and a grown-up can write one more idea.
\newcommand{\feedbackcard}{%
%% alt: A feedback card that says Your card! with five dashed sticker spots next to the dance, the sound, the surprise, the story and the colors, and a line for one more idea.
  \begin{tikzpicture}[spot/.style={circle,draw=kcgray,dashed,minimum size=0.8cm,inner sep=0pt}]
    \draw[kcpurple,line width=1.2pt,rounded corners=6pt] (0,0) rectangle (8.2,8.0);
    \node[anchor=north west,font=\bfseries] at (0.3,7.8) {\faGift\ Your card!};
    \node[anchor=north west,font=\small] at (0.3,7.1) {Put a sticker by what you liked.};
    \foreach \icon/\word [count=\i] in {\faRedo/the dance,\faVolumeUp/the sound,\faMousePointer/the surprise,\faBookOpen/the story,\faPalette/the colors}
      {\node[spot] at (0.9,6.7-\i*1.0) {};
       \node[anchor=west] at (1.6,6.7-\i*1.0) {\icon\ \word};}
    \node[anchor=north west,font=\small] at (0.3,1.05) {One more idea:};
    \draw[kcgray!60] (0.3,0.3) -- (7.9,0.3);
  \end{tikzpicture}}
\begin{center}
\feedbackcard\hfill\feedbackcard\par\bigskip
\feedbackcard\hfill\feedbackcard
\end{center}
\end{document}
